% Notebook 17 module. No preamble. Unique label prefix cs-.
% Requires the usual norm macro and jt-real-theorem from the joint-tail module.
\section{The Gaussian at an arbitrary current renormalization scale}
\label{cs-section}

The fixed-ratio bare Gaussian and the Gaussian used after earlier
classical minimizations are different objects. We now extract the
uniform bounds for the latter. They are already contained in
Notebooks 05--07; this section supplies the current-scale probability
normalization and records how Notebook 15 removes the earlier
homogeneous smooth-minimizer hypothesis.

\subsection{The retained coordinates and their exact precision}

Set the current lattice spacing to one. The original spacing is
$M^{-1}$, with $M=L^k$, while the next blocking ratio $L\ge3$ is fixed
and odd. All regulators are compatible periodic tori. Let
\[
 q_M(B)=\inf_{Q_M^{\rm st}A=B}\norm{d_{1/M}A}_{1/M}^2
\]
be the classical flat effective quadratic form. The fine norm has
volume factor $M^{-d}$ and the current norm is counting measure.
Next-step maps are denoted by $Q_L^\flat,T_L,K_L$: respectively the
rooted constraint, averaged root-to-site contour map, and pivot
insertion. These are the fixed-$L$ maps of Notebook 05. The map
$K_L$ inserts a field in $\ker Q_L^\flat$ and is the identity on
the retained nonpivot coordinates. It is not a gauge-selected
fine minimizing extension.

The actual flat exponential-gauge quadratic model on the retained
variables $z$ is
\begin{equation}\label{cs-precision}
 A_{M,0}=K_L^T(q_M+T_L^TT_L)K_L,\qquad
 d\mu_{M,0}(z)=
 (2\pi)^{-n/2}\det(A_{M,0})^{1/2}
 e^{-z^TA_{M,0}z/2}\,dz .
\end{equation}
Here $n$ counts independent real bond/color coordinates once, using
a real color basis orthonormal for the normalized trace. Thus
\begin{equation}\label{cs-counting}
 \operatorname{Cov}_{\rm count}(z)=A_{M,0}^{-1}.
\end{equation}
There is no factor $M^d$ in \eqref{cs-counting}: these are
\emph{current-unit} coordinates, not original fine potential values.
The coupling has already been removed from the Gaussian by the
change of field $B=g_kK_Lz$ at linear order. The nonlinear constraint
correction, Haar density and coordinate Jacobian are kept outside
the chosen Gaussian in \eqref{cs-precision}. In particular, a
coupling-suppressed quadratic Haar term may be part of the
interaction; no vanishing of that term is asserted.

\subsection{Uniform spectral bounds, including all previous depths}

The normalized cochain contraction in Notebook 06 gives
\begin{equation}\label{cs-energy-sandwich}
 \norm{d_1B}_1^2\le q_M(B)\le C_d\norm B_1^2,
 \qquad C_d=2^{d+6}(d-1).
\end{equation}
The compact lift proves the upper bound for odd $M\ge3$; $M=1$
uses the identity lift and satisfies the same displayed bound.
These estimates hold for all current fields, not just a transverse
Fourier sector.

For clarity we record the fixed geometric constants of Notebook 06.
On the cube of $3^d$ neighboring $L$-blocks form the finite matrix
\[
 \mathcal M_L=d_\Omega^Td_\Omega+
 T_{0,\Omega}^TT_{0,\Omega}+
 Q_{{\rm st},\Omega}^TQ_{{\rm st},\Omega},
 \quad \sigma(d,L)=\lambda_{\min}(\mathcal M_L)>0.
\]
Its kernel is trivial: zero curl gives a gradient, the tree rows
make its scalar potential constant on each block, and the averaging
rows equate neighboring block constants. Summing the local
inequality gives the axial gap
$\gamma_{\rm loc}=\sigma(d,L)/3^d$ on
$\ker Q_L^{\rm st}\cap\ker T_0$. Periodic repetitions are counted
with multiplicity, so the same constant applies on short compatible
tori.

Put
\begin{equation}\label{cs-constants}
 \begin{split}
 d_0&=2\sqrt d,\qquad
 t_0=\sqrt{d(L-1)(L^d-1)},\qquad
 r_L=1+2d(L-1)/L,\\
 K_*^2&=1+L^{2d}(2d+2)r_L^2,\\
 \gamma_*&=
 \left((1+d_0t_0)^2/\gamma_{\rm loc}+d_0^2\right)^{-1},
 \qquad B_*=(C_d+t_0^2)K_*^2.
 \end{split}
\end{equation}
The flat gauge bridge of Notebook 05 gives
\[
 q_M(B)+\norm{T_LB}^2\ge\gamma_*\norm B^2
 \quad (Q_L^\flat B=0).
\]
One can see the constraint compatibility directly. First subtract
the rooted gradient $D(T_LB)$ to obtain $B_T$ with
$T_LB_T=Q_L^{\rm st}B_T=0$. Convert it by a block-mean-zero
gradient to the tree representative $B_0$. Their curls agree and
$\norm{B_T}\le(1+d_0t_0)\norm{B_0}$. Apply the axial gap and
weighted Cauchy--Schwarz to
$B=B_T+D(T_LB)$. This argument uses the lower inequality in
\eqref{cs-energy-sandwich}; it does not confuse the rooted and
straight constraints.

Since $z$ is a coordinate restriction of $K_Lz$ and
$\norm{K_L}\le K_*$, the full precision sandwich is
\begin{equation}\label{cs-spectral}
 \gamma_*I\le A_{M,0}\le B_*I,
 \qquad
 B_*^{-1}I\le A_{M,0}^{-1}\le\gamma_*^{-1}I.
\end{equation}
All constants are independent of $k$, $M$ and the torus volume.
There is no need to replace this proved bound by an assumed
transverse ellipticity statement.

\subsection{Weighted locality and its scope}

Group the retained coordinates by their fixed next-step $L$-cells.
For a block matrix $F$, use the weighted Schur norm
\[
 \norm F_\kappa=
 \max\left\{
 \sup_i\sum_j e^{\kappa|i-j|_1}\norm{F_{ij}},
 \sup_j\sum_i e^{\kappa|i-j|_1}\norm{F_{ij}}
 \right\},
\]
with periodic distance on finite regulators. Notebook 07,
its flat-decay theorem and section on periodic gluing and the reduced inverse,
prove constants $\kappa_*>0,C_*<\infty$, depending only on fixed
geometry, such that
\begin{equation}\label{cs-weighted}
 \norm{A_{M,0}}_{\kappa_*}
 +\norm{A_{M,0}^{-1}}_{\kappa_*}\le C_*.
\end{equation}
The proof there first establishes a uniform holomorphic Fourier
strip for $q_M$, then uses fixed-range $K_L,T_L$, the real gap
\eqref{cs-spectral}, and a smaller strip for the inverse.
Exact periodization gives the same conclusion on compatible tori.
This additional locality does not follow from a spectral sandwich
alone. It concerns the flat retained-coordinate matrix; it is not
strict locality in a varying background.

\subsection{Joint value-and-difference tails on the current lattice}

Let $\nabla_1$ contain all forward unit differences of every bond
component, and define
\begin{equation}\label{cs-tested-field}
 Y=\begin{pmatrix}K_Lz\\\nabla_1K_Lz\end{pmatrix}.
\end{equation}
The elementary bound
$\norm{\nabla_1}\le2\sqrt d$ and
\eqref{cs-spectral} imply
\begin{equation}\label{cs-tested-covariance}
 0\le\operatorname{Cov}_{\rm count}(Y)\le
 K_{\rm cur}I,
 \qquad K_{\rm cur}=(1+4d)K_*^2/\gamma_*.
\end{equation}
Because $K_L$ and the difference stencil have fixed range,
\eqref{cs-weighted} also supplies exponential block decay for this
tested covariance, with a fixed loss in its constant.

If each tested cell uses at most $m$ scalar coordinates and every
coordinate is tested in at most $\nu$ cells, the real joint-tail
theorem~\ref{jt-real-theorem} gives
\begin{equation}\label{cs-joint-tail}
 \mathbb P\left(
 \max_{i\in I_b}|Y_i|>R\text{ for every }b\in S\right)
 \le
 \left[m\sqrt2\exp\left\{-\frac{R^2}{4\nu K_{\rm cur}}\right\}
 \right]^{|S|}.
\end{equation}
Fixed current-scale cells, including fixed collars for derivatives,
have $m,\nu$ independent of $M$. The Gaussian is correlated;
independence is neither asserted nor needed. For $g_kY$, replace
$R$ by $R/g_k$. If $R(g)=c(\log g^{-2})^p$ with $p>1/2$, the
bracket at radius $R(g)$ is smaller than any fixed power of $g$
below a sufficiently small ceiling.

To restore the physical current spacing $a$, write
$Z_{\rm phys}(x)=a^{-1}(K_Lz)(x/a)$ in four dimensions. Then
\[
 (aZ_{\rm phys},a^2\nabla_aZ_{\rm phys})=Y.
\]
This is the current-step normalization. It differs from treating
all original fine variables between $M^{-1}$ and one as a single
fluctuation. It also differs from the bare fixed-ratio formula
$(LI,\nabla_{1/L})C_L(LI,\nabla_{1/L})^T$, whose coordinates use
next-coarse spacing one. Both statements are consistent after their
scales and Gaussian models are specified.

\subsection{What is now identified with the deposited source}

Notebook 07's source section records the chain from the deposited
small-field step to \eqref{cs-precision}: Raw II:73865--73950
straightens the constraint, rescales by $g_k$, and eliminates the
pivots; Raw I:46960--47000 fixes the pivot coefficient convention.
At the trivial field the action Hessian is $q_M$, the exponential
gauge Hessian is $T_L^TT_L$, and the action gradient vanishes.
Consequently the quadratic start of the constraint-straightening
correction does not change this Hessian. Haar and Jacobian factors
remain in the interaction.

The smooth-minimizer premise in that earlier identification is now
supplied, in the homogeneous setting, by Notebook 15's theorem on
transfer to the ordinary homogeneous averaging constraint. It
constructs an analytic local branch for the ordinary iterated group
average, with real action minimality in a fixed small global
potential class. At each finite regulator the ordinary constraint
is an analytic submersion at the identity. Its derivative is onto,
as follows from the fixed-step pivot right inverses and their
composition. Thus every tangent-feasible field is the tangent to
a feasible curve. Comparing the local minimum with these curves
shows that its second derivative is the quadratic variational
infimum for the rooted derivative. Notebook 06's explicit
gradient adjustment identifies that infimum with the displayed
$Q_M^{\rm st}$ infimum. The gradient adjustment preserves energy;
it does not identify the two averaging maps as maps on potentials.

Accordingly \eqref{cs-precision} is the homogeneous flat source
Gaussian under the deposited ordinary-average and normalization
conventions already matched in those notebooks. The nonlinear
branch need not equal an independently chosen axial representative:
the locally minimizing action and its flat Hessian agree modulo
gauge. This does not re-certify every original journal convention,
extend the result to mixed histories or holes, or identify a
nonflat independent-current analytic family.

\subsection{A distinct fixed-slice perturbation consequence}

There is also a useful quadratic model without pivot coordinates.
Let $X_L=\ker Q_L^{\rm st}\cap\ker(R_LD_1^T)$ be the fixed flat
source-gauge slice for the unit-to-$L$ step, and let $J_L$ be its
isometric inclusion. The fixed-$L$ Maxwell covariance proved in
Notebooks 09 and 12 implies
\[
 J_L^Td_1^Td_1J_L\ge c_LI
\]
for some $c_L>0$, uniformly in compatible volume. For example one
may take $c_L=(L^2\norm{C_L})^{-1}$, where $C_L$ uses next-coarse
spacing one. Hence
\begin{equation}\label{cs-slice-gap}
 c_LI\le J_L^Tq_MJ_L\le C_dI.
\end{equation}
This gives the same kind of current-scale value-and-difference
tail bound on $X_L$, uniformly in $M$.

Notebook 16 proves that the ordinary coarse classical Hessian
satisfies $D^2V_M(B)=q_M+O(\norm B_\infty)$ in weighted row and
column norms. Its compression to this fixed slice therefore has a
relative perturbation bounded by $C\norm B_\infty/c_L$. This is
a valid input to the abstract complex relative-quadratic tail
theorem on a sufficiently small common background ball. It is an
effective classical Gaussian family in a fixed linear slice.
It is not automatically the source's nonlinear fiber Gaussian:
the latter also contains chart-curvature/current terms, exponential
gauge fixing, and an exact cutoff-coordinate prescription.
Notebook 19 now performs that additional physical-branch
identification, keeping both chart-curvature terms and the
actual absolute gauge function.
Nor is weighted locality of this compressed inverse inferred by
conjugating with weights which do not preserve $X_L$; the weighted
inverse assertion \eqref{cs-weighted} concerns the separately
proved pivot-coordinate family.
